% =====================================================================
\chapter{Что-то не работает?}
\label{ch:trouble}
% =====================================================================
\chapterintro
  {как по шагам найти, где «теряется» сигнал; частые проблемы и их решения}
  {терпение и эта глава}

\section{Ищем, где потерялся сигнал}

\begin{simple}
Сигнал идёт по цепочке: стик → EdgeTX → радиомодуль → приёмник → полётный контроллер → моторы.
Если модель не слушается, надо найти \textbf{первое звено}, где сигнал пропал. Иди по цепочке
с начала и на каждом шаге задавай вопрос «да или нет?».
\end{simple}

\begin{figure}[H]
\centering
\begin{tikzpicture}[
  q/.style={draw=etx,fill=etx!6,rounded corners=3pt,text width=62mm,align=center,minimum height=10mm,font=\sffamily\small},
  a/.style={draw=warn,fill=warn!6,rounded corners=3pt,text width=52mm,align=left,font=\sffamily\scriptsize},
  ok/.style={draw=okgreen,fill=okgreen!8,rounded corners=3pt,text width=62mm,align=center,font=\sffamily\small},
  yes/.style={font=\sffamily\scriptsize\bfseries,text=okgreen!50!black},
  no/.style={font=\sffamily\scriptsize\bfseries,text=warn}]
  \node[draw=black,fill=black!8,rounded corners=3pt,font=\sffamily\small\bfseries] (s) at (0,0) {Модель не слушается стиков};
  \node[q] (q1) at (0,-1.5) {1. На странице каналов пульта полоски двигаются от стиков?};
  \node[a,anchor=west] (a1) at (4.0,-1.5) {калибровка (гл.~\ref{ch:radio}); выключенная строка во входе или миксере; выбрана не та модель};
  \node[q] (q2) at (0,-3.3) {2. Светодиод приёмника горит постоянно? В скрипте ELRS растёт счётчик пакетов?};
  \node[a,anchor=west] (a2) at (4.0,-3.3) {привязка: фраза, версия ELRS, \menu{Internal RF} = \menu{CRSF}, Model Match (гл.~\ref{ch:rf})};
  \node[q] (q3) at (0,-5.1) {3. В Betaflight на вкладке Receiver полоски двигаются?};
  \node[a,anchor=west] (a3) at (4.0,-5.1) {провода RX--FC (TX↔RX), UART и протокол CRSF в Betaflight};
  \node[q] (q4) at (0,-6.9) {4. Двигаются \emph{правильные} полоски (газ = Throttle)?};
  \node[a,anchor=west] (a4) at (4.0,-6.9) {\menu{Channel map} = AETR1234 в Betaflight и порядок каналов в пульте};
  \node[q] (q5) at (0,-8.7) {5. Коптер армится?};
  \node[a,anchor=west] (a5) at (4.0,-8.7) {диапазон ARM на вкладке Modes; газ не на нуле; в CLI команда \code{status} покажет причину};
  \node[ok] (e) at (0,-10.3) {Всё работает --- можно лететь!};
  \draw[arr] (s) -- (q1);
  \foreach \i [evaluate=\i as \j using int(\i+1)] in {1,2,3,4}{\draw[arr] (q\i) -- node[right,yes]{да} (q\j);}
  \draw[arr] (q5) -- node[right,yes]{да} (e);
  \foreach \i in {1,...,5}{\draw[arr,warn] (q\i.east) -- node[above,no]{нет} (a\i.west);}
\end{tikzpicture}
\caption{Пошаговый поиск неисправности}
\end{figure}

\section{Частые проблемы}

\begin{small}
\renewcommand{\arraystretch}{1.25}
\begin{longtable}{@{}L{50mm}L{110mm}@{}}
\toprule
\textbf{Что случилось} & \textbf{Почему и что делать} \\
\midrule
\endhead
Пульт не включается & аккумуляторы сели или вставлены неправильно; держи кнопку питания
  дольше; у 18650 проверь, достают ли они до контактов \\
При включении «SD card warning» & версия файлов на SD-карте не та: обнови SD через Buddy
  (гл.~\ref{ch:firmware}) \\
Застрял на «Throttle / Switch warning» & опусти газ, переведи переключатели в стартовое
  положение; если требуемое положение невозможно --- исправь \menu{Preflight checks} \\
Стик в центре не 0, по краям не 100 & повтори калибровку (гл.~\ref{ch:radio}) \\
Среднее положение тумблера не работает & в \menu{Hardware} стоит \menu{2POS} --- поставь \menu{3POS} \\
В скрипте ExpressLRS «\texttt{-{}-}» или «Loading…» & \menu{Internal RF} выключен или не \menu{CRSF};
  попробуй \menu{Baudrate} 400K \\
ELRS не привязывается & разные фразы; разные основные версии ELRS; разная частота (2,4\,ГГц и
  900\,МГц); включён Model Match с другим номером \\
FrSky не привязывается & не совпадает подтип (D8/D16/LBT/FCC) или версия (v1/v2);
  попробуй \menu{RF Freq. fine tune} \\
Газом крутится коптер & перепутан порядок каналов: AETR в пульте и \code{AETR1234} в Betaflight \\
Коптер не армится & \sw{CH5} не попадает в диапазон ARM; газ не на нуле; коптер стоит криво;
  в CLI \code{status} покажет причину \\
Нет телеметрии & \menu{Telem Ratio} = \menu{Off}; повтори \menu{Discover new sensors} \\
Связь рвётся рядом с пультом & повреждена антенна приёмника или она закрыта карбоном; на
  большой мощности вплотную приёмник «глохнет» --- отойди на 2--3\,м \\
Пульт молчит & \menu{Sound mode} без голоса; громкость 0; нет папки \code{SOUNDS/ru} или в
  \menu{Voice language} другой язык \\
Компьютер не видит пульт & кабель «только для зарядки»; не выбран режим USB; для прошивки на
  Windows --- драйвер WinUSB (Zadig) \\
Lua-скрипт не запускается & не та папка, длинное имя файла, ошибка в коде \\
Boxer: вентилятор шумит всё время & большая мощность ELRS; уменьши мощность или подними \menu{Fan Thresh} \\
Быстро садится батарея пульта & мощность 500--1000\,мВт, яркая подсветка; включи динамическую мощность \\
\bottomrule
\end{longtable}
\end{small}

\section{Хорошие привычки}

\begin{itemize}
  \item Раз в месяц --- резервная копия SD-карты.
  \item Перед сезоном --- проверка калибровки и failsafe.
  \item Одна настоящая модель --- одна модель в пульте; подпиши имя и на самой модели.
  \item Заряжай пульт перед каждым выездом, а не «вроде ещё хватит».
  \item Любые проверки моторов --- \textbf{только без пропеллеров}.
\end{itemize}
